% 2015-05-21 - Emerson Ribeiro de Mello - mello@ifsc.edu.br
% \documentclass[handout,xcolor=pdftex,dvipsnames,table]{beamer}
\documentclass{beamer}

\usepackage[utf8]{inputenc}
\usepackage[T1]{fontenc}
\usepackage[portuguese]{babel}

% usando tema personalizado. 
% arquivo beamerthemeIFSC.sty deve estar no mesmo diretório do .tex
\usepackage{beamerthemeIFSC}
\usepackage{tikz}
\usepackage{xargs}
\usepackage{amsmath}
\usepackage{mathtools}
\usepackage{soul} %fazer texto tachado
\usepackage{epigraph}
\usepackage[colorinlistoftodos,prependcaption,textsize=tiny]{todonotes}
\usepackage{graphicx}
\usepackage{listings}
\usepackage{color}

\definecolor{dkgreen}{rgb}{0,0.6,0}
\definecolor{gray}{rgb}{0.5,0.5,0.5}
\definecolor{mauve}{rgb}{0.58,0,0.82}

\lstset{frame=tb,
	language=Java,
	aboveskip=3mm,
	belowskip=3mm,
	showstringspaces=false,
	columns=flexible,
	basicstyle={\small\ttfamily},
	numbers=none,
	numberstyle=\tiny\color{brown},
	keywordstyle=\color{blue},
	commentstyle=\color{dkgreen},
	stringstyle=\color{mauve},
	breaklines=true,
	breakatwhitespace=true,
	tabsize=3
}

\usepackage{xargs}
%
\newcommandx{\unsure}[2][1=]{\todo[linecolor=red,backgroundcolor=red!25,bordercolor=red,#1]{#2}}
\newcommandx{\info}[2][1=]{\todo[linecolor=OliveGreen,backgroundcolor=OliveGreen!25,bordercolor=OliveGreen,#1]{#2}}
%

\hypersetup{pdfstartview={Fit},pdftitle={\@title},
 	pdfsubject={IFBA},pdfauthor={\@author}
}

\hyphenation{ob-je-ti-vo}


%%%%%%%%%%%%%%%%%%%%%%%%%%%%%%%%%%%%%%%%%%%%


\title{Python}
\subtitle{Programando de forma idiomática}
\author{Eliakin Costa}
\date{16 de junho de 2018}
\institute{Instituto Federal da Bahia\\
campus Salvador\\
\url{eliakin.costa@kde.con}
}



%%%%%%%%%%%%%%%%%%%%%%%%%%%%%%%%%%%%%%%%%%%%

\begin{document}

\begin{frame}[t]
	\maketitle
\end{frame}

% Descomente as linhas abaixo se desejar colocar um sumário de todas as seções
\begin{frame}[t]{Agenda}
\tableofcontents
\end{frame}

% O trecho de código abaixo serve para sobrescrever regras do template
\def\sectionname{}
\def\insertsectionnumber{}
\def\subsectionname{}
\def\insertsubsectionnumber{}

\AtBeginSection{\frame{\sectionpage}\addtocounter{framenumber}{-1}}


\AtBeginSubsection{\frame{\subsectionpage}\addtocounter{framenumber}{-1} }
\AtBeginSubsubsection{\frame{\subsubsectionpage}\addtocounter{framenumber}{-1} }



%%%%%%%%%%%%%%%%%%%%%%%%%%%%%%%%%%%%%%%%%%%%
% Inicio do documento
%%%%%%%%%%%%%%%%%%%%%%%%%%%%%%%%%%%%%%%%%%%%

\begin{frame}{Quem sou eu?}
    \begin{itemize}
        \item Meu nome é Eliakin Costa,
        \item Sou aluno de ADS (IFBA),
        \item Apaixonado pelo software livre $\backslash$o/
        \item Colaborador do KDE
        \item \textbf{Contatos}:
        \begin{itemize}
            \item \textbf{E-mail}: eliakin.costa@kde.com
            \item \textbf{GitHub}: github.com/eliakincosta
        \end{itemize}
    \end{itemize}
\end{frame}
%

\section{Introdução}
\begin{frame}{O começo de tudo}
    \begin{itemize}
        \item  CWI - Centro Tecnológico
        \item Guido Van Rossum
        \item Linguagem ABC        
    \end{itemize}
\end{frame}
%
\begin{frame}{Python!? Que nome é esse? }
	\begin{center}
		\includegraphics[scale=0.5]{img/python.png}
	\end{center}
\end{frame}

\begin{frame}{Python!? Que nome é esse? }
    \begin{itemize}
        \item Python Flying in Circus
        \item Episódio do spam xD
		\item Guido não gostava da relação do python com uma cobra        
    \end{itemize}
\end{frame}
%

\begin{frame}{Python!? Que nome é esse? }
\begin{center}
	\includegraphics[scale=1.0]{img/python_book.jpg}
\end{center}
\end{frame}

\begin{frame}{Guido Van Rossum - O ditador benevolente}
\begin{itemize}
	\item Está vivo!!!!
	\item Google de 2005-2013
	\item Dropbox de 2013-hoje
\end{itemize}
\end{frame}

\begin{frame}{Por que eu devo aprender Python?}
\begin{itemize}
        \item Linguagem de Propósito geral
        \item Simples e intuitivo
        \item Comunidade espetacular
        \item Você consegue fazer quase tudo apenas com os recursos da  própria linguagem
        \item Multiparadigma
    \end{itemize}
\end{frame}
%

\begin{frame}{Quem está usando python?}
\begin{figure}[htp]
	\centering
	\includegraphics[width=.3\textwidth]{img/bit_torrent.jpg}\quad
	\includegraphics[width=.3\textwidth]{img/globo.png}\quad
	
	\medskip
	\medskip
	\medskip
	\includegraphics[width=.3\textwidth]{img/google.png}\quad
	\includegraphics[width=.3\textwidth]{img/youtube.png}
	
	\medskip
\end{figure}

\end{frame}
%

\begin{frame}{Quem está usando python?}
\begin{figure}[htp]
	\centering
	\includegraphics[width=.3\textwidth]{img/blender.jpg}\quad
	\includegraphics[width=.3\textwidth]{img/gimp.png}\quad
	\includegraphics[width=.3\textwidth]{img/krita.png}\quad
\end{figure}

\end{frame}


\section{Tipagem Dinâmica, mas nem tanto}
\begin{frame}{Características da Tipagem Dinâmica}

\begin{itemize}
	\item Verificações de tipo feitas em runtime
	\item Verificações feita sobre o valor(ex: class int)
	\item Bob Harper, um dos criadores da linguagem Standard ML, "uma linguagem dinamicamente tipada é uma linguagem estaticamente tipada com apenas um tipo estático".
\end{itemize}
\end{frame}

\begin{frame}{Como é a tipagem estática?}

\centering \text{JAVA}
\lstset{language=Java}
\lstinputlisting{code/tipagem_estatica.java}
      
\end{frame}

\begin{frame}{Como é a tipagem dinâmica?}

\centering \text{Python}
\lstset{language=Python}
\lstinputlisting{code/tipagem_dinamica.py}

\end{frame}

\begin{frame}{Tipagem dinâmica, mas forte}
\begin{itemize}
	\item Observe o seguinte código em php:
	\lstset{language=php}
	\lstinputlisting{code/tipagem_fraca.php}
\end{itemize}
\begin{itemize}
	\item Vamos testar o mesmo código em python.
\end{itemize}

\end{frame}

\begin{frame}{Tipagem forte}
	\begin{itemize}
		\item Código mais robusto
		\item Tipos inválidos ocasionarão exceções
		\item Explícito é melhor do que implícito
	\end{itemize}
\end{frame}


\section{Tipos, Instruções e Estilo}
\begin{frame}{Tipos Básicos do python}
\begin{itemize}
	\item Númericos:
	\begin{itemize}
		\item inteiro (int)
		\item ponto flutuante (float)
		\item booleano (bool) 1 ou 0
		\item complexo (complex)
	\end{itemize}
\end{itemize}
\begin{itemize}
		\item Obs:
		\begin{itemize}
			\item Suportam as operações básicas, mod, etc (+, -, *, /, \%, **, +=, -=, *=, /=,\%=, **=).
			\item Operações entre diferentes tipo irão resultar no tipo mais complexo
			\item Cuidado, números inteiros tem precisão infinita
		\end{itemize}
\end{itemize}
\end{frame}

\begin{frame}{Tipos Básicos do python}
\begin{itemize}
	\item Iteráveis:
	\begin{table}[]
		\centering
		\caption{Iteráveis}
		\label{my-label}
		\begin{tabular}{|l|c|c|l|l|}
			\hline
			\multicolumn{1}{|c|}{} & Ordenada              & Modificável           & \multicolumn{1}{c|}{Unicidade}  \\ \hline
			strings                   & \multicolumn{0}{c|}{X} &                           &                       \\ \hline
			listas                 & X                     & X                     &                                          \\ \hline
			tuplas                 & X                     & \multicolumn{1}{l|}{} &                    \\ \hline
			sets                   & \multicolumn{1}{l|}{} & X                     & \multicolumn{1}{c|}{X}                             \\ \hline
		\end{tabular}
	\end{table}
	\item Vamos tentar um pouco de código
\end{itemize}
		\textbf{\text{Obs: Iremos focar apenas em listas}}
\end{frame}

\section{Funções e Classes}
\begin{frame}{Funções}

\begin{itemize}
	\item São definidas pela palavra reservada \textit{def}
	\item Podem possuir parâmetros default
	\item São objetos de primeira classe
	\begin{itemize}
		\item Criada em tempo de execução
		\item Pode ser atribuída a uma variável
		\item Passada como argumento para uma função		
		\item Devolvida como resultado de uma função
	\end{itemize}
\end{itemize}
\end{frame}


\begin{frame}{Funções}

\centering \text{Exemplo de função}
\lstset{language=Python}
\lstinputlisting{code/exemplo_funcao.py}
\end{frame}

\begin{frame}{Classes}

\begin{itemize}	
	\item São definidas com a palavra reservada \textit{class}
	\item Python suporta herança múltipla
	\item Python possui sobrecarga de operadores, assim como C++
	\item Todos os atributos são públicos em Python
	\item Python não possui uma palavra reservada para interface como JAVA, mas o python tem classes abstratas.
	
\end{itemize}

\end{frame}

\begin{frame}{Classes }
\centering \text{Exemplo de classe}
\lstset{language=Python}
\lstinputlisting[
basicstyle=\tiny, %or \small or \footnotesize etc.
]{code/exemplo_classe.py}
\end{frame}

\section{Prática}
\begin{frame}{Praticando!}
    \begin{itemize}
        \item Comentários
        \item Macros "$\%$TODO"
        \item Indentação
        \item Separação de sílabas
        \item Iniciando nova linha
        \item Separação de arquivos
        \item Customizando texto
        \begin{itemize}
            \item \textit{itálico}
            \item \textbf{negrito}
            \item \underline{sublinhado}
            \item \textst{tachado}\footnote{Usa pacote específico}
        \end{itemize}
     \end{itemize}
\end{frame}

\begin{frame}{Praticando! Mergulhando fundo...}
    \begin{itemize}
        \item Nota de Rodapé
        \item Criando Epígrafes:
            \begin{itemize}
                \item \textit{epigraph}\footnote{Usa pacote específico}
                \item \textit{quote}
                \item \textit{description}
            \end{itemize}
        \item Criando Listas:
            \begin{itemize}
                \item \textit{itemize}
                \item \textit{enumerate}
                \item \textit{description}
            \end{itemize}
                
    \end{itemize}
\end{frame}

\begin{frame}{Praticando! Ao infinito e Além!}
    \begin{itemize}
        \item Criando Recursos Gráficos:
        \begin{itemize}
            \item Imagem
            \item Tabela
            \item Fórmulas
        \end{itemize}
        \item Criando Citações
        \begin{itemize}
            \item Seção
            \item Imagem
            \item Tabela
        \end{itemize}
        \item Criando novos comandos
        \item Criando referências 
        \item Citação de referências
        \item Alterando \textit{layouts}
        
    \end{itemize}
\end{frame} 


%Finalizando
\begin{frame}{Peguntas}
	\includegraphics[scale=0.5]{img/perguntas.jpg}
\end{frame} 
 
\begin{frame}{Referências}
 	\begin{itemize}
 		\item Python Fluente - Luciano Ramalho
 		\item Learning Python - Mark Lutz
 		\item Python Tutorial https://docs.python.org/3/tutorial/
 	\end{itemize}
\end{frame} 


\end{document}